\setcounter{zagc}{6}\zag{Funkcje procesowe i funkcje przepisów prawa karnego procesowego}[Ł]
\noindent\textbf{Funkcja procesowa}\medskip

Funkcja procesowa jest to zasadniczy typ działalności procesowej wyznaczony przez procesową rolę odpowiedniego podmiotu procesu karnego (uczestników postępowania karnego)\medskip

\textbf{Funkcja procesowa ścigania} – zmierzanie do ustalenia, ujęcia i pociągnięcia do odpowiedzialności karnej sprawcy przestępstwa, przy czym należy pamiętać, że funkcja ścigania nie pokrywa się w pełni z funkcją oskarżania gdyż ściganie realizują nie tylko oskarżyciele, lecz także organy prowadzące postępowanie przygotowawcze, nadzorujące oraz wspomagające to postępowanie, np. policja wykonująca zlecone czynności śledcze.\medskip

\textbf{Ściganie karne} to działalność procesowa skierowana na wykrycie i w razie potrzeby ujęcie sprawcy przestępstwa oraz pociągnięcie go do odpowiedzialności karnej.\begin{itemize}
    \item W sprawach publicznoskargowych ściganie rozpoczyna się z chwilą pierwszej czynności organu postępowania przygotowawczego skierowanej na wykrycie lub ujęcie sprawcy;
    \item W sprawach prywatnoskargowych ściganie rozpoczyna się wraz z pierwszą czynnością procesową oskarżyciela prywatnego zmierzająca do pociągnięcia oskarżonego do odpowiedzialności.
\end{itemize}

\textbf{Stadium postępowania przygotowawczego}\begin{itemize}
    \item Podmioty ścigania karnego – organy śledcze i dochodzeniowe, a także organy nadzoru;
    \item Działalność procesowa organów postępowania przygotowawczego skierowana jest na:\begin{itemize}
        \item Ustalenie, czy zostało popełnione przestępstwo,
        \item Wykrycie i w razie potrzeby ujęcie sprawcy przestępstwa,
        \item Wyjaśnienie okoliczności sprawy, w tym ustalenie osób pokrzywdzonych i rozmiarów szkody,
        \item Zebranie, zabezpieczenie i w niezbędnym zakresie utrwalenie dowodów;
    \end{itemize}
    \item Realizacja zasady obiektywizmu (art. 4 \acr{kpk}) oraz swobodnej oceny dowodów (art. 7 \acr{kpk});
    \item Podmiotem realizującym funkcję ścigania karnego w stadium przygotowawczym jest także pokrzywdzony, któremu prawo karne procesowe przyznaje w tym stadium uprawnienia strony.
\end{itemize}

\textbf{Stadium postępowania jurysdykcyjnego} – ściganie karne przybiera formę oskarżenia:\begin{itemize}
    \item Funkcję tę spełnia oskarżyciel – publiczny, posiłkowy, prywatny;
    \item W postępowaniu sądowym oskarżyciel jest stroną zasadniczo pozbawioną uprawnień władczych.
\end{itemize}

\textbf{Stadium wykonawcze} – ściganie przyjmuje postać zarządzenia wykonania kary, które należy do sądu pierwszej instancji oraz w postaci nadzoru nad wykonaniem kary.\medskip

\textbf{Funkcja obrony karnej} – reakcja na ściganie; ochrona interesów oskarżonego w procesie. Funkcję obrony mogą spełniać wszystkie osoby mające prawo działania w procesie w interesie oskarżonego, a więc:\begin{itemize}
    \item Sam oskarżony;
    \item Jego obrońca;
    \item Oskarżyciel publiczny oraz w pewnym zakresie;
    \item Osoba najbliższa dla oskarżonego, jego przedstawiciel ustawowy lub osoba, pod której opieką oskarżony faktycznie pozostaje.
\end{itemize}

Oskarżony, jego obrońca i prokurator są uprawnieni zasadniczo do podejmowania wszelkich czynności obrończych, z wyjątkiem nielicznych wypadków, w których jest wymagane osobiste działanie oskarżonego lub prokuratora bądź obrońcy. Natomiast oskarżyciele publiczni niebędący prokuratorami oraz osoby najbliższe są uprawnione do obrony tylko w odniesieniu do niektórych kategorii oskarżonych (np. nieletni, ubezwłasnowolnieni) lub niektórych tylko czynności obrończych.\medskip

\textbf{Funkcja orzekania (rozstrzygania)} – rozpoznanie sprawy w procesie karnym i rozstrzygnięcie o jej przedmiocie, czyli kwestii odpowiedzialności prawnej (karnej i ewentualnie cywilnej) oskarżonego. Funkcja rozstrzygania nie jest też tożsama z funkcją orzekania, gdyż decyzje procesowe rozstrzygające jakieś kwestie w procesie karnym podejmują nie tylko sądy, ale także prokurator, np. o wszczęciu postępowania lub o wstąpieniu do już wszczętego procesu. Rozpoznanie sprawy obejmuje:\begin{itemize}
    \item Przeprowadzenie i ocenę dowodów oraz ustalenie na ich podstawie faktów;
    \item Zastosowanie do ustalonych faktów właściwej normy karnoprawnej – przyjęcie prawidłowej kwalifikacji prawnej czynu;
    \item Wyciągnięcie prawnych konsekwencji z ustalonych faktów.
\end{itemize}

Rozpoznanie $\neq$ wyjaśnienie i rozstrzyganie sprawy!\begin{itemize}
    \item Ten typ działalności występuje wszędzie tam, gdzie zachodzi konieczność wydania decyzji procesowej. Dlatego funkcja ta odnosi się także do postępowania przygotowawczego, w którym również są podejmowane decyzje procesowe.
\end{itemize}

Podmiotami funkcji rozstrzygania sprawy są:\begin{enumerate}
    \item W stadium przygotowawczym – organy śledztwa i dochodzenia oraz organy nadzoru. Są one w tym stadium równocześnie podmiotami ścigania karnego;
    \item W stadium jurysdykcyjnym – sąd, prezes sądu (przewodniczący wydziału, sędzia upoważniony) oraz przewodniczący składu orzekającego;
    \item W stadium wykonawczym – podmioty wymienione w art. 2 \acr{kkw}, organy postępowania wykonawczego oraz organy postępowań incydentalnych w zakresie, w jakim wydają decyzje, opinie lub występują z wnioskami.
\end{enumerate}

\textbf{Odpowiedzialność cywilna oskarżonego}:\begin{itemize}
    \item Tożsamość funkcji procesowych – ściganie, obrona i rozstrzyganie;
    \item Konieczną przesłankę stanowi odpowiedzialność karna oskarżonego;
    \item Podmiotem dochodzenia roszczeń cywilnych w procesie karnym jest pokrzywdzony lub podmiot wykonujący prawa pokrzywdzonego. 
\end{itemize}

\noindent\textbf{Funkcje przepisów prawa karnego procesowego}\medskip

Cel to pewien postulowany stan rzeczy, który ma być osiągnięty dzięki podjęciu określonych czynności. Cel jest czymś zamierzonym, zaplanowanym przez tego, kto podejmuje działanie. Funkcja oznacza rzeczywisty, obiektywny skutek, rezultat działania jakiejś instytucji lub normy postępowania; jeżeli zaplanowany skutek wystąpił, wówczas pojęcie funkcji pokrywa się z celem. Wyróżnia się następujące funkcje przepisów prawa karnego procesowego:\medskip

\textbf{Funkcja prakseologiczna (instrumentalna)} – ukształtowanie procesu karnego w sposób pozwalający na osiągnięcie jego celów, a zatem:\begin{itemize}
    \item Sprawiedliwości materialnej – pociągnięcie do odpowiedzialności jedynie sprawcy przestępstwa, któremu udowodniono winę w sposób przewidziany prawem; oraz
    \item Sprawiedliwości proceduralnej – ukształtowanie procesu karnego w sposób pozwalający na sprawiedliwe rozstrzygnięcie o jego przedmiocie po przeprowadzeniu rzetelnego postępowania karnego.
\end{itemize}

\textbf{Funkcja gwarancyjna} – ochrona uczestników procesu karnego przed niezgodną z prawem działalnością organów procesowych, czyli podjętą bez podstawy prawnej i poza granicami prawa. Funkcja gwarancyjna chroni też przed bezprawną działalnością podjętą przez innych uczestników postępowania, a także świadków i biegłych (ekspertów) i inne osoby wykonujące w procesie wyznaczone im przez prawo role.

\textbf{Funkcja porządkująca} – przepisy prawa karnego procesowego wyznaczają porządek czynności, ich sekwencję, spełniają rolę koordynatora czynności procesowych. Ustalają prawa i obowiązki organów oraz innych uczestników postępowania, wyznaczają im swoiste, konwencjonalne role. Stwarzają prawną podstawę, granice działania i określenie formy podejmowanych czynności.
